% !TEX root = ../ADEmain.tex
%\def\manifold{\operatorname{ma}}
\def\manifold{\mathcal{M}}
\def\subM{{_{\! \manifold}}}

\def\numb{N}
\def\numbd{\numb^{\circ}}
%\def\nmin{\numb_{\operatorname{loc}}}
\def\nmin{\numb_{l}}
\def\nlin{\numb_{\operatorname{lin}}}
\def\nman{\numb_{_{\! \manifold}}}

\def\Kern{\mathcal{K}}
\def\sigmaX{\sigma_{\!_{X}}}
\def\Xbar{\Mv}
\def\UUV{\mathcal{V}}
\def\corr{\tau}

\def\LL{\mathcal{L}}

\def\jd{j^{\circ}}
\def\jjl{j\jl}
\def\jl{\alpha}
\def\jld{\jl^{\circ}}

%\def\weightM{\omega}
\def\weightM{\delta}
\def\WeightM{\Delta}

\def\svd{\sv^{\circ}}
\def\xvm{\xv^{\dag}}
\def\displ{\nu}

\def\JJJ{\mathcal{J}}
%\def\nJJJ{\numb_{\!_{\JJJ}}}
\def\nJJJ{J}
\def\JJJd{\JJJ^{\circ}}
\def\init{\operatorname{in}}
\def\outp{\operatorname{out}}
\def\IEDR{\mathcal{A}}

\def\fcns{\fcn^{*}}
\def\fcls{\fcl^{*}}
\def\fclvs{\fclv^{*}}

\def\sumI{\sum_{i \in \II}}


%%%%%%%%%%%%%%%%%%%%%%%%%%%%%%%  new section  %%%%%%%%%%%%%%%%%%%%%%%%%%%%%%%
\Section{A single-index model}
\label{Ssingind}

Consider the e.d.r. regression problem
\begin{EQA}
	Y_{i}
	&=&
	\fs(\betav^{\T} \Xv_{i}) + \eps_{i} \, ,
\label{YifXieiedrSI}
\end{EQA}
where \( \Xv_{i} \in \R^{\dimd} \) for large \( \dimd \), \( \betav \) is an index vector and
\( \fs \) is a univariate \emph{link function}.
Introduce also a function \( \gs(x) = \E \bigl( Y \cond \Xv=\xv \bigr) \).
A structural single-index assumption yields the relations
\begin{EQA}
	\gs(\xv)
	&=&
	\fs(\betav^{\T} \xv) ,
	\\
	\nabla \gs(\xv)
	&=&
	\fs'(\betav^{\T} \xv) \betav.
\label{fxgKxedrSI}
\end{EQA} 
In particular, the derivative of \( \gs \) at any point \( \xv \) is proportional
to the target vector \( \betav \).
This suggests recovering \( \betav \) from the nonparametrically estimated gradient of \( \gs \).


%%%%%%%%%%%%%%%%%%%% section %%%%%%%%%%%%%%%%%%%%
\Subsection{Preliminaries}
%Define \( \Psi(\xv) = \bigl( \psi_{m}(\xv) \bigr) \in \R^{\dimd+1} \) for \( \psi_{0}(\xv) \equiv 1 \), 
%\(	\psi_{m}(\xv) =	\langle \xv ,\ev_{m} \rangle \), \( m = 1,\ldots,\dimd \). 
Suppose a subset \( \II \) of the whole index set \( \{ 1,\ldots,n \} \) to be given.
We refer to \( \II \) as a training set.
A default choice is \( \II = \{ 1,\ldots,n \} \).
For a collection of centers \( \{ \xv_{j} \}_{j \in \JJJ} \) with \( |\JJJ| = \nJJJ \),
localizing tensor \( \TEDR \), 
and a localizing kernel \( \Kern \), consider localizing weights
\begin{EQA}
	\weight_{ij}
	&=&
	\Kern\bigl( \| \TEDR(\Xv_{i} - \xv_{j}) \|^{2} \bigr) 
	\, ,
	\qquad
	i \in \II,
	\quad
	j \in \JJJ
	\, .
\label{s7cuqkdwywdhbjwbfjwbdj}
\end{EQA}
At the first step, we apply \( \TEDR = \hiso^{-2} \Id_{\dimd} \)
leading to isotropic weights 
\( \weight_{ij} = \Kern\bigl( \hiso^{-2} \| \Xv_{i} - \xv_{j} \|^{2} \bigr) \),
where the bandwidth \( \hiso \) has to ensure sufficiently many points in almost every vicinity 
\( U_{\hiso}(\xv_{j}) \) of \( \xv_{j} \) of radius \( \hiso \). 
For each \( j \in \JJJ \), define the local mass and mean
\begin{EQA}[c]
	\neff_{j}
	\eqdef
	\sumI \weight_{ij}
	\, ,
	\qquad
	\Xbar_{j}
	=
	\frac{1}{\neff_{j}}
	\sumI \Xv_{i} \, \weight_{ij}
%	\, ,
%	\qquad
%	j \in \JJJ
	\, .
\label{uyhvvy883hierwfewAVE}
\end{EQA}
%The main difference with the approach of Section~\ref{Ssingind} is the use of directional averaging. 
For each point \( \xv_{j} \), suppose a set \( \Sph_{j} \) of directions (unit vectors) \( \sph \in \R^{\dimd} \) 
to be given.
Compute for every \( \sph \in \Sph_{j} \)
\begin{EQ}[rcl]
	\Ima_{j,\sph}
	& \eqdef &
	\frac{1}{\neff_{j}}
	\sumI Y_{i} \, \langle \Xv_{i} - \Xbar_{j} \, , \sph \rangle \, \weight_{ij}
	\, ,
	\\
	\Uv_{j,\sph}
	& \eqdef &
	\frac{1}{\neff_{j}}
	\sumI (\Xv_{i} - \Xbar_{j}) \, \langle \Xv_{i} - \Xbar_{j} \, ,\sph \rangle  \, \weight_{ij}
	\in \R^{\dimd}
	\, .
	\qquad
\label{juu8hjfioee433erhAVE}
\end{EQ}
The model equation \eqref{YifXieiedrSI} yields
\begin{EQA}[c]
	\E \Ima_{j,\sph}
	=
	\frac{1}{\neff_{j}}
	\sumI \fs(\betav^{\T} \Xv_{i}) \, \langle \Xv_{i} - \Xbar_{j} \, ,\sph \rangle \, \weight_{ij}
	\, .
\end{EQA}
Moreover, an approximation \( \fs(t) \approx \fcn_{j} + \fcl_{j} (t - t_{j}) \) implies by definition of \( \Xbar_{j} \)
\begin{EQA}[rcl]
	\E \Ima_{j,\sph}
	& \approx &
	\sumI \bigl\{ \fcn_{j} + \fcl_{j} \, \betav^{\T} (\Xv_{i} - \xv_{j}) \bigr\} \, \langle \Xv_{i} - \Xbar_{j} \, ,\sph \rangle \, \weight_{ij}
	\\
	&=&
	\frac{1}{\neff_{j}}
	\sumI \fcl_{j} \, \betav^{\T} (\Xv_{i} - \Xbar_{j}) \, \langle \Xv_{i} - \Xbar_{j} \, ,\sph \rangle \, \weight_{ij}
%	=
%	\fcl_{j} \, \uv_{j}^{\T} \KSme_{j} \betav
	=
	\fcl_{j} \, \Uv_{j,\sph}^{\T} \, \betav
	\, .
\end{EQA}
%where
%\begin{EQA}[c]
%	\Uv_{j,\sph}
%	\eqdef
%	\sumI (\Xv_{i} - \Sv_{j}) \, \langle \Xv_{i} - \Sv_{j} \, ,\sph \rangle  \, \weight_{ij}
%	\, .
%	\qquad \quad
%\label{uyhvvy883hierwfewAVE}
%\end{EQA}
%Now we define
%\begin{EQA}
%	\LL_{\lambda}(\prmtv)
%	&=&
%	\frac{1}{2} \| \Imav - \imav \|^{2} 
%	+ \frac{1}{2} \sum_{j \in \JJJ} \Bigl\| \imav_{j} - \fcl_{j} \, \Uv_{j}^{\T} \betav \Bigr\|^{2} 
%	+ \frac{\lambda}{2} \| \betav - \betav_{\init} \|^{2} \, .
%\label{tjjc7c6f5fAVE}
%\end{EQA}
With \( \Imav_{j} = (\Ima_{j,\sph})_{\sph \in \Sph_{j}} \in \R^{\nSph} \) and a \( \nSph \times \dimd \)-matrix 
\( \UV_{j} = (\Uv_{j,\sph})^{\T} \), this leads to the minimization of
\begin{EQA}[c]
	\sum_{j \in \JJJ} \sum_{\sph \in \Sph_{j}} \bigl( \Ima_{j,\sph} - \fcl_{j} \, \Uv_{j,\sph}^{\T} \, \betav \bigr)^{2} 
%	+ \lambda \| \betav - \betav_{\init} \|^{2}
	=
	\sum_{j \in \JJJ} \bigl\| \Imav_{j} - \fcl_{j} \, \UV_{j} \, \betav \bigr\|^{2}
%	+ \lambda \| \betav - \betav_{\init} \|^{2}
	\qquad
%	\to
%	\inf_{\betav,(\fcl_{j})_{j \in \JJJ}}
%	\, ,
\label{ikdfjoiefjoieuAVE}
\end{EQA}
over all unit vectors \( \betav \in \R^{\dimd} \) and \( (\fcl_{j})_{j \in \JJJ} \in \R^{\nSph} \). 
%which can be solved by alternating minimization as for \eqref{ikdfjoiefjoieuvbjmr3e}.
One can also underweight the terms \( j \) with a relatively small value \( \neff_{j} \).
This leads to minimization of 
\begin{EQA}[c]
	\sum_{j \in \JJJ} \bigl\| \Imav_{j} - \fcl_{j} \, \UV_{j} \, \betav \bigr\|^{2} \,
	\neff_{j}
	\, .
\end{EQA}

\Subsubsection{Prediction and generalization from the single-index model}
With \( \betav \) and \( \fcl_{j} \) fixed, estimation of the values \( \fcn_{j} \approx \gs(\xv_{j}) \) is straightforward.
Consider the sums
\begin{EQA}[c]
	S_{j}
	\eqdef
	\frac{1}{\neff_{j}}
	\sumI Y_{i} \, \weight_{ij}
	=
	\frac{1}{\neff_{j}}
	\sumI \gs(\Xv_{j}) \, \weight_{ij}
	+  
	\frac{1}{\neff_{j}}
	\sumI \eps_{i} \, \weight_{ij}
	\, .
\end{EQA}
Model equation \( \gs(\xv) = \fs(\xv^{\T} \betavs) \) and
a local approximation \( \fs(t) \approx \fcns_{j} + \fcls_{j} (t - t_{j}) \) imply by definition \eqref{uyhvvy883hierwfewAVE}
of \( \Xbar_{j} \)
\begin{EQA}[c]
	\E S_{j}
	\approx
	\fcn_{j} + \frac{\fcls_{j}}{\neff_{j}} \sumI (\Xv_{i} - \xv_{j})^{\T} \betavs \; \weight_{ij}
	=
	\fcns_{j} + \fcls_{j} (\Xbar_{j} - \xv_{j})^{\T} \betavs 
	\, .
\end{EQA}
This suggests estimating the value \( \fcns_{j} = \gs(\xv_{j}) \) by
\begin{EQA}[c]
	\fcn_{j}
	=
	S_{j} - \fcl_{j} (\Xbar_{j} - \xv_{j})^{\T} \betav
	\, .
\end{EQA}
This leads to the fit
\begin{EQA}[rcl]
	\gs(\xv)
	& \approx &
	\fcn_{j} + \fcl_{j} (\xv - \xv_{j})^{\T} \betav
	=
	S_{j} + \fcl_{j} (\xv - \Xbar_{j})^{\T} \betav
\end{EQA}
in a vicinity of the point \( \xv_{j} \).
In the case of a regular design, the value \( |(\xv_{j} - \Xbar_{j})^{\T} \betav| \) is small and one can use 
a simplified rule
\begin{EQA}[rcl]
	\gs(\xv_{j})
	& \approx &
	S_{j}
	=
	\frac{1}{\neff_{j}}
	\sumI \gs(\Xv_{j}) \, \weight_{ij}
	\, .
\label{jfuyujudeyudy7dyuew}
\end{EQA}

\Subsection{Data fit}
The proposed procedure is iterative and involves several steps.
A proper stopping rule is an issue and requires evaluating the quality of estimation at every step.
%
The quality of estimation will be measured by the generalization error on the testing set
which corresponds to the set \( \JJJ \) of centers \( \xv_{j} \).
One can apply a simple twofold approach by splitting the data into a training set \( \II \) and a testing set \( \JJJ \).
The training set \( (Y_{i},\Xv_{i})_{i \in \II} \) will be used for estimating the parameters of the procedure.
%
For each \( j \in \JJJ \), assume that the point \( \Xv_{i(j)} \) is aligned  
with a center \( \xv_{j} \).
The use of \eqref{jfuyujudeyudy7dyuew} leads to the fit
\begin{EQA}[rcl]
	\err
	& \eqdef &
%	\sum_{j \in \JJJ} \bigl\{ Y_{i(j)} - \fcn_{j} - \fcl_{j} (\Xv_{i} - \xv_{j})^{\T} \betav \bigr\}^{2} 
%	\\
%	&=&
	\sum_{j \in \JJJ} \bigl\{ Y_{i(j)} - S_{j} \bigr\}^{2} 
\label{jvfujkiiiut5tyerjfd}
\end{EQA}
which can be used, e.g., for a data-driven stopping rule.
As the structural adaptation step excludes an overfitting effect, one can use the whole dataset for training:
\( \II \eqdef \{ 1,\ldots,n \} \).
At the initialization step, we use another representation for \( \err \); see Section~\ref{Siniloclin}.


%Let a point \( \xv \) be close to some \( \xv_{j} \) for \( j \in \JJJ \).
%Then \( \gs(\xv) \)can be approximated by 
%\begin{EQA}[c]
%	\fcn_{j} + \fcl_{j} (\xv - \xv_{j})^{\T} \betav 
%	=
%	S_{j} + \fcl_{j} (\xv - \Xbar_{j})^{\T} \betav
%	\, .
%\end{EQA}
%A more flexible approach averages over all \( j \in \JJJ \) with the corresponding weights 
%\( \weight_{j} = \Kern(\| \TEDR (\xv - \xv_{j}) \|^{2}) \):
%\begin{EQA}[c]
%	\gs(\xv)
%	\approx
%	\biggl\{ \sum_{j \in \JJJ} \Kern(\| \TEDR (\xv - \xv_{j}) \|^{2}) \biggr\}^{-1}
%	\sum_{j \in \JJJ} \bigl\{ S_{j} + \fcl_{j} (\xv - \Xbar_{j})^{\T} \betav \bigr\} 
%		\, \Kern(\| \TEDR (\xv - \xv_{j}) \|^{2}) 
%\end{EQA}

\Subsection{Alternating optimization (AO)}
\label{SLLaltern}
This section discusses an approach to solving the minimization problem \eqref{ikdfjoiefjoieuAVE}.
It assumes that the vectors
\( \Imav_{j} = (\Ima_{j,\sph})_{\sph \in \Sph_{j}} \in \R^{\nSph} \) and \( \nSph \times \dimd \)-matrices
\( \UV_{j} = (\Uv_{j,\sph})^{\T} \) for \( j \in \JJJ \) are available.
The main challenge is that the unit sphere \( \Sphere_{\dimd} \) is not convex in \( \R^{\dimd} \), 
and minimization of \eqref{ikdfjoiefjoieuAVE} over this set is numerically hard.


\paragraph{Relaxation and regularization}
We apply the usual relaxation trick by allowing at every step of the minimization procedure 
an unconstrained optimization over \( \betav \in \R^{\dimd} \) and then by projecting 
the obtained solution to the unit sphere.
Additional numerical stability can be ensured 
by a penalty \( \mfrac{\lambda}{2} \| \betav - \betav_{\init} \|^{2} \) with 
\( \betav_{\init} \in \R^{\dimd} \)
meaning a starting guess or the previous step estimate.
%One can use \( \betav_{\init} = (1,0,\ldots,0)^{\T} \) or a random unit vector un \( \R^{\dimd} \).
%It is only important that \( \langle \betav,\ev_{1} \rangle \) significantly deviates from zero.
With a tuning parameter \( \lambda \), this leads to an iterative procedure which minimizes at each step 
the objective function
\begin{EQA}
	\LL_{\lambda}(\prmtv)
	&=&
	\frac{1}{2} \sum_{j \in \JJJ} \bigl\| \Imav_{j} - \fcl_{j} \, \UV_{j} \betav \bigr\|^{2} \,	
	\neff_{j}
	+ \frac{\lambda}{2} \| \betav - \betav_{\init} \|^{2} 
	\, .
\label{tjjc7c6f5fd5e5erjwyhfd}
\end{EQA}
%Now we sketch the alternating optimization (AO) procedure for minimization of \eqref{tjjc7c6f5fd5e5erjwyhfd}.
%Also, a starting guess \( \betav_{\init} \) is required.
%
\emph{One-step alternating} reads as follows:

\begin{mylist}
\item
For \( \betav_{\init} \) fixed and each \( j \in \JJJ \), the value \( \fcl_{j} \) is estimated by
\begin{EQA}[c]
	\hat{\fcl}_{j}
	=
	\frac{\langle \Imav_{j}, \UV_{j} \, \betav_{\init} \rangle}{\| \UV_{j} \, \betav_{\init} \|^{2}} 
	\, .
\end{EQA}
%where \( \tau \geq 0 \) is a regularization parameter which prevents dividing by small values of \( \| \UV_{j} \, \betav_{\init} \| \).

\item
For given \( (\hat{\fcl}_{j})_{j \in \JJJ} \), the vector \( \betav \) can be estimated as
\begin{EQA}[c]
	\hat{\betav}
	=
	\biggl( 
		\sum_{j \in \JJJ} \hat{\fcl}_{j}^{\; 2} \; \UV_{j}^{\T} \UV_{j} \,
		\neff_{j} 
		+ \lambda \Id_{\dimd} 
	\biggr)^{-1}
	\biggl( 
		\sum_{j \in \JJJ} \hat{\fcl}_{j} \, \UV_{j}^{\T} \Imav_{j} \, \neff_{j}
		+ \lambda \betav_{\init} 
	\biggr)
	\, .
\end{EQA}
{\color{blue} please check.}

\item
The output \( \betav_{\outp} \) is obtained by normalization of \( \hat{\betav} \):
\begin{EQA}[c]
	\betav_{\outp}
	=
	\frac{1}{\| \hat{\betav} \|} \hat{\betav}
	\, .
\end{EQA}

\end{mylist}

\paragraph{Alternating optimization procedure} % of \eqref{tjjc7c6f5fd5e5erjwyhfd}} 
repeats the one-step alternating several times.
At every iteration, the output \( \betav_{\outp} \) of the previous step is used as input \( \betav_{\init} \).
The procedure delivers the unit vector \( \betav_{\outp} \) after \( k_{\max} \) iterations.



%%%%%%%%%%%%%%%%%%%% section %%%%%%%%%%%%%%%%%%%%
\Subsection{Anisotropic weights and structural adaptation}
\label{SanisoSI}
The local weights \( \weight_{ij} \) from \eqref{s7cuqkdwywdhbjwbfjwbdj} involve a local tensor \( \TEDR \).
The single-index model structure \( \gs(\xv) = \fs(\betav^{\T} \xv) \) in \eqref{YifXieiedrSI} suggests
anisotropic weights
\begin{EQA}
	\weights_{ij}
	&=&
	\Kern\bigl( h^{-2} |(\Xv_{i} - \xv_{j})^{\T} \betavs|^{2} \bigr) .
\label{s7cuqkdwywdhbjwbfjwbdji}
\end{EQA}
Unfortunately, such defined weights require knowledge of the index vector \( \betavs \) and 
a correctly specified single-index model. 
We follow the structure-adaptive approach from \cite{HJPS}, which applies anisotropic weights 
with an increasing anisotropy parameter during iterations.
As mentioned earlier, for initialization, we apply a non-informative choice 
\begin{EQA}[c]
	\weight_{ij} = \Kern\bigl( \hiso_{0}^{-2} \| \Xv_{i} - \xv_{j} \|^{2} \bigr) 
	\, .
\end{EQA}
The bandwidth \( \hiso_{0} \) can be selected as the smallest value ensuring
\begin{EQA}[c]
	\frac{1}{\nJJJ}\sum_{j \in \JJJ} \sumI \Kern\Bigl( \mfrac{\| \Xv_{i} - \xv_{j} \|^{2}}{\hiso_{0}^{2}} \Bigr)
	\geq 
	\nmin
	\, .
\label{jhgiverkytde4w34etycj}
\end{EQA}
This condition means that in a local vicinity of each point \( \xv_{j} \) there are (on average) 
at least \( \nmin \) design points \( \Xv_{i} \).

Anisotropic weights assume some prior information about the structure of the model, which is accumulated 
in the localizing tensor \( \TEDR \).
Namely, for two values \( \hiso \) and \( \aniso \leq 1 \) 
and a current guess \( \betav_{\init} \), anisotropic weights \( \weight_{ij} \)
correspond to the tensor 
\begin{EQA}[rcl]
	\TEDR^{2}
	&=&
	\hiso^{-2} \bigl\{ \aniso^{2} \, (\Id_{\dimd} - \betav_{\init} \betav_{\init}^{\T}) 
	+ \betav_{\init} \betav_{\init}^{\T} \bigr\}
	=
	\hiso^{-2} \aniso^{2} \, \Id_{\dimd} 
	+ (1 - \aniso^{2}) \hiso^{-2} \betav_{\init} \betav_{\init}^{\T} 
	\, 
\label{iufhvu8kjbmnll475edf}
\end{EQA}
leading to 
\begin{EQA}[c]
	\weight_{ij}
	=
	\Kern\bigl( \| \TEDR (\Xv_{i} - \xv_{j}) \|^{2} \bigr)
	=
	\Kern\Bigl( 
		\mfrac{\aniso^{2} \| \Xv_{i} - \xv_{j} \|^{2}}{\hiso^{2}} 
		+ \mfrac{(1 - \aniso^{2}) \langle \Xv_{i} - \xv_{j}, \betav_{\init} \rangle^{2}}{\hiso^{2}} 
	\Bigr)
	\, .
\end{EQA}
The quantity \( \aniso \) describes the anisotropy strength.
Small values of \( \hiso, \aniso \) correspond to strong localization in direction \( \betav = \betav_{\init} \)
and mild localization in orthogonal directions. 
During iteration, the procedure decreases \( \hiso \) and \( \aniso \) in a way 
to ensure sufficiently many design points in the local vicinity of \( \xv_{j} \) 
for almost every \( j \in \JJJ \).
For fixed \( \hiso \) and \( \betav \), the anisotropy parameter \( \aniso \) is the largest value ensuring
\begin{EQA}[c]
	\frac{1}{\nJJJ} \sum_{j \in \JJJ} 
	\sumI \Kern\bigl( \| \TEDR (\Xv_{i} - \xv_{j}) \|^{2} \bigr)
%	=
%	\frac{1}{\nJJJ} \sum_{j \in \JJJ} \sumI \Kern\Bigl( 
%		\mfrac{\aniso^{2} \| \Xv_{i} - \xv_{j} \|^{2}}{\hiso^{2}} + \mfrac{\langle \Xv_{i} - \xv_{j}, \betav \rangle^{2}}{\hiso^{2}} 
%	\Bigr)
%	=
%	\frac{1}{\nJJJ} \sum_{j \in \JJJ} \sumI \weight_{ij}
	\geq 
	\nmin
	\, .
	\qquad
\label{jhgiverkytde4w34etycja}
\end{EQA}
It is easy to see that a decay of \( \hiso \) leads to an increase of \( \hiso/\aniso \).

%%%%%%%%%%%%%%%%%%%% section %%%%%%%%%%%%%%%%%%%%
\Subsection{Design of experiments}
\label{SdesignSI}
Fix the values \( d \), \( n \), \( \sigma_{\eps} \), \( \sigmaX \) such that \( n/\sigma_{\eps}^{2} \geq 20 d \).
Finally, fix a unit vector \( \betavs \) and a univariate function \( \fs(x) \).
%; e.g \( f(x) = \sin(s x) \) or \( f(x) = x \sin(s x) \) for \( s=1,2,3,4 \).
\begin{mylist}

\item 
generate a correlated design \( \Xv_{i} = \sigmaX \bigl\{ \corr \gaussv_{0} + (1 - \corr) \gaussv_{i} \bigr\} \),
where \( \gaussv_{i} \) are i.i.d. standard Gaussian in \( \R^{\dimd} \), \( i=0,1,\ldots,n \);
\item 
generate errors \( \eps_{i} \sim \ND(0,\sigma_{\eps}^{2}) \) and compute \( Y_{i} = \fs(\Xv_{i}^{\T} \betavs) + \eps_{i} \);

%\item
%generate \( \betav_{\ini} \) uniformly on the unit sphere in \( \R^{\dimd} \).
\end{mylist}

{\color{blue} 
\begin{mylist}
\item
provide possibility to select \( n \in \{ 800, \bm{1000}, 1200, 2000 \} \), \( \dimd \in \{ 10, 20,\ldots, 90, \bm{100} \} \), 
\( \sigma_{\eps} \in \{ 0, 0.1, \bm{0.2}, 0.4, 0.6, 0.8, 1 \} \), \( \corr \in \{ 0,0.2,\bm{0.4},0.8 \} \), \( \fs(x) \).
Include \( \fs(x) = \sin(sx) \), \( \bm{f(x)} = x \sin(s x) \) for \( s=1,2,\bm{3},4 \).
(Default choices in boldface.)
\item

Plot the data \( Y_{i} \) against \( \Xv_{i}^{\T} \betavs \).
\item
Preview the possibility of re-drawing the whole dataset.
\end{mylist}
}

%%%%%%%%%%%%%%%%%%%% section %%%%%%%%%%%%%%%%%%%%
\Subsection{Initialization by local linear estimation}
\label{Siniloclin}
The objective function \( \LL(\prmtv) \) is not convex. 
Therefore, the starting direction \( \betav_{\init} \) is important for a sensitive final result.
\cite{HJPS} suggested fixing \( \betav_{\init} \)
by local linear estimation of the gradient of \( \gs \) at the points \( \xv_{j} \).
Namely, fix \( \hiso_{\lin} \) as the smallest value ensuring 
\begin{EQA}[c]
	\frac{1}{\nJJJ}\sum_{j \in \JJJ} 
	\sumI \Kern\Bigl( \mfrac{\| \Xv_{i} - \xv_{j} \|^{2}}{\hiso_{\lin}^{2}} \Bigr)
	\geq 
	\nlin
	\, .
\label{hhfiecredgvhihiue86t}
\end{EQA}
%The default value of \( \nlin \) is \( n/\nmin \).
%The proposal requires \( \nlin > \dimd + 1 \).
Now, for each \( j \in \JJJ \), compute
\begin{EQA}[rcl]
	\neff_{j}
	\eqdef
	\sumI \Kern\Bigl( \mfrac{\| \Xv_{i} - \xv_{j} \|^{2}}{\hiso_{\lin}^{2}} \Bigr)
	\, ,
	\qquad
	\Xbar_{j}
	& \eqdef &
%	\biggl\{ \sumI \Kern\Bigl( \mfrac{\| \Xv_{i} - \xv_{j} \|^{2}}{\hiso_{\lin}^{2}} \Bigr) \biggr\}^{-1}
	\frac{1}{\neff_{j}}
	\sumI \Xv_{i} \, \Kern\Bigl( \mfrac{\| \Xv_{i} - \xv_{j} \|^{2}}{\hiso_{\lin}^{2}} \Bigr)
	\, ,
	\qquad
%	\\
%	S_{j}
%	& \eqdef &
%	\biggl\{ \sumI \Kern\Bigl( \mfrac{\| \Xv_{i} - \xv_{j} \|^{2}}{\hiso_{\lin}^{2}} \Bigr) \biggr\}^{-1}
%	\sumI Y_{i} \, \Kern\Bigl( \mfrac{\| \Xv_{i} - \xv_{j} \|^{2}}{\hiso_{\lin}^{2}} \Bigr)
%	\, ,
%	\qquad
%	\\
%	\UUV_{j}
%	& \eqdef &
%	\sumI \binom{1}{\Xv_{i} - \Xbar_{j}} \, \binom{1}{\Xv_{i} - \Xbar_{j}}^{\T}
%		\Kern\Bigl( \mfrac{\| \Xv_{i} - \xv_{j} \|^{2}}{\hiso_{\lin}^{2}} \Bigr)
	\, ,
\end{EQA}
and
\begin{EQA}[rcl]
	\UUV_{j}
	& \eqdef &
	\frac{1}{\neff_{j}}
	\sumI (\Xv_{i} - \Xbar_{j}) \, (\Xv_{i} - \Xbar_{j})^{\T}
		\Kern\Bigl( \mfrac{\| \Xv_{i} - \xv_{j} \|^{2}}{\hiso_{\lin}^{2}} \Bigr)
	\, ,
	\\
	S_{j}
	& \eqdef &
	\frac{1}{\neff_{j}}
	\sumI Y_{i} (\Xv_{i} - \Xbar_{j}) \, 
		\Kern\Bigl( \mfrac{\| \Xv_{i} - \xv_{j} \|^{2}}{\hiso_{\lin}^{2}} \Bigr)
	\, .
\end{EQA}
%Given a regularization coefficient \( \tau \), 
Define
\begin{EQA}[rcl]
%	\binom{\hat{\gs}(\xv_{j})}{\hat{\nabla \gs}(\xv_{j})}
	\hat{\nabla \gs}(\xv_{j})
	&=&
%	(\UUV_{j} + \tau \Id_{\dimd})^{-1} S_{j}
	\UUV_{j}^{-1} S_{j}
	\, .
\label{ihvifvityguwhd8354e}
\end{EQA}
Define \( \betav_{\init} \) as the first principal direction
for the set of the estimated gradients \( \{ \hat{\nabla \gs}(\xv_{j}) \}_{j \in \JJJ} \).
For this, compute
\begin{EQA}[c]
	\JEDR
	\eqdef
	\sum_{j \in \JJJ} \hat{\nabla \gs}(\xv_{j}) \{ \hat{\nabla \gs}(\xv_{j}) \}^{\T}
	\neff_{j}
	\, .
\label{uhcudiiy66yhjikfed}
\end{EQA}
Then \( \betav_{\init} \) is the eigenvector of \( \JEDR \) corresponding to its largest eigenvalue
\( \lambda_{1}(\JEDR) \).
If \( \lambda_{1}(\JEDR) > \lambda_{2}(\JEDR) \) then this eigenvector is uniquely defined up to directional sign.
However, we only need the product \( \betav_{\init} \, \betav_{\init}^{\T} \) which is uniquely defined.

%\paragraph{Prediction quality \( \err_{\init} \)} of the initialization step is measured as in \eqref{jvfujkiiiut5tyerjfd}:
%\begin{EQ}[rcl]
%	\err_{\init}
%	& \eqdef &
%	\sum_{j \in \JJJ} 
%	\bigl\{ Y_{i} - S_{j(i)} - \hat{\fcl}_{j(i)} (\Xv_{i} - \Xbar_{j(i)})^{\T} \betav_{\init} \bigr\}^{2} 
%	\, ,
%\label{jvfujkiiiut5tyerjfdIni}
%	\\
%	\hat{\fcl}_{j}
%	&=&
%	\frac{\langle \Imav_{j}, \UUV_{j} \, \betav_{\init} \rangle}{\| \UUV_{j} \, \betav_{\init} \|^{2}}
%	\, ,
%	\qquad
%	j \in \JJJ
%	\, .
%\end{EQ}
%\begin{EQ}[rcl]
%	\err_{\init}
%	& \eqdef &
%	\sum_{j \in \JJJ} 
%	\bigl\{ Y_{i(j)} - \hat{\gs}(\xv_{j}) \bigr\}^{2} 
%	\, .
%\label{jvfujkiiiut5tyerjfdIni}
%\end{EQ}


%%%%%%%%%%%%%%%%%%%% section %%%%%%%%%%%%%%%%%%%%
\Subsection{Main procedure}

This section describes the structure-adaptive procedure.
The procedure consists of two big parts.
The first step is initialization; it mainly produces the starting fit \( \betav_{\init} \).
Then the procedure iterates the structure-adaptive estimation step with a decreasing anisotrophy parameter
\( \aniso_{k} \).

%%%%%%%%%%%%%%%%%%%% section %%%%%%%%%%%%%%%%%%%%
\Subsubsection{Main procedure. Initialization}
This step follows the idea from Section~\ref{Siniloclin}.
Tuning parameters include
\( \nJJJ \), \( \nlin \), \( k_{\max} \), \( \displ \); see Table~\ref{TdefSIini}.
\begin{table}[h]
\begin{tabular}{c|c|c}
        \hline
        Parameter & Meaning & Possible values and default choice \\
        \hline
        \( \nlin \) & local mass for local linear estimation & 
        \( d + s \nmin \) for \( s= 1,2,\bm{3} \) \\
        \( \nJJJ \) & size of \( \JJJ \) & \( s \, n/\nmin \) for \( s=1, \ldots , \bm{\nmin/2},\ldots,\nmin \) \\
        \( \displ \)  & displacement coefficient & 0, \textbf{0.1}, 0.5, 1 \\
%        \( \tau \)  & regularization parameter & 0, 0.05, \textbf{0.1}, 0.5, 1 \\
		\( \Kern \) & the kernel & Epanechnikov: \( \Kern(t) = (1-t^{2})_{+} \) \\
		\hline
    \end{tabular}
\caption{Tuning parameters for initialization}    
\label{TdefSIini}
\end{table}

The procedure reads as follows.
\begin{mylist}

\item 
\textbf{Fix a collection of points \( (\xv_{j})_{j \in \JJJ} \).}
Given \( \nJJJ \), select randomly indices \( i(1), \ldots, i(\nJJJ) \) from \( \{ 1,\ldots,n \} \).
Set \( \xv_{j} = \Xv_{i(j)} + \displ \, \sigmaX \gaussv_{j} \) with \( \gaussv_{j} \sim \ND(0,1) \)
and a displacement coefficient \( \displ \) for \( j \in \JJJ \).
The observations \( (Y_{i(j)})_{j \in \JJJ} \) compose the \emph{testing set}.

\item
\textbf{The training set \( \II \)} is defined either by as \textbf{the whole collection} \( \bm{\II = \{ 1,\ldots,n \} } \)
or as a collection of the remaining indices \( i \notin \{ i(1), \ldots, i(\nJJJ) \} \).
%For each \( i = i(j) \in \II_{\test} \), define \( j(i) = j \), so that \( \Xv_{i} \) is aligned with \( \xv_{j(i)} \). 
%\\
%{\color{blue} Provide both options.}

\item
\textbf{Fix} \( \betav_{\init} \) using local linear regression and PCA as
in Section~\ref{Siniloclin}.

%\item
%Compute the testing error \( \err_{\init} \) by \eqref{jvfujkiiiut5tyerjfdIni}.
%\begin{EQA}[rcl]
%	\err_{\init}
%	&=&
%	\sum_{j \in \JJJ} \bigl\{ Y_{i} - S_{j(i)} - \fcl_{j(i)} (\Xv_{i} - \Xbar_{j(i)})^{\T} \betav_{\init} \bigr\}^{2} 
%	\, .
%\label{jvfujkiiiut5tyerjfdI}
%\end{EQA}

\item
Fix \( \hiso_{0} \) as the smallest value ensuring \eqref{jhgiverkytde4w34etycj}.

\item
Set \( k=0 \).
Define \( \TEDR_{0} = \hiso_{0}^{-1} \Id_{\dimd} \).
\end{mylist}

{\color{blue} 
\begin{mylist}
	\item
	Provide a possibility to manually select the tuning parameters \( \nlin \), \( \nJJJ \).

	\item
	Display \( \bigl| \cos\langle \betavs,\betav_{\init} \rangle \bigr| \).
	
	\item
	Display the first and the second eigenvalues of \( \JEDR \).
	
	\item
	Provide a possibility to optimize \( \nlin \)
	to get the best initialization results.
	
%	\item
%	{\color{blue} Explore the impact of \( \nJJJ \).}
	
\end{mylist}
}

%%%%%%%%%%%%%%%%%%%% section %%%%%%%%%%%%%%%%%%%%
\Subsubsection{Main procedure. Structural adaptation.}
For the procedure, we have to fix a set of tuning parameters; see Table~\ref{TdefSI}
for the complete list and default values.
%
Parameters \( \nJJJ \) and \( \nSph \) are mainly for controlling the numerical
complexity of the procedure.
Parameter \( \lambda \) has the meaning of the step size, 
it controls stability and convergence of alternating optimization
from Section~\ref{SLLaltern}.
Parameter \( \nmin \) is essential for the structure-adaptive block.
However, the procedure exhibits a very stable performance 
for a minor variation of the parameters around the default values (in boldface).
%
\begin{table}[h]
\begin{tabular}{c|c|c}
        \hline
        Parameter & Meaning & possible values and default \\
        \hline
        \( \nmin \) & \( \# \) of design points \( \Xv_{i} \) in a local vicinity & \( 7, 10, 15, \bb{20} \) \\
%        \( \nJJJ \) & size of \( \JJJ \) 				& \( s \, n/\nmin \) for \( s=1, \ldots , \bm{\nmin/2},\ldots,\nmin \) \\
		\( \nSph \) & size of each \( \Sph_{j} \)		& \( \bm{\nmin}, \, 2\nmin, \, 3\nmin \) \\
%\( a (= 2^{1/\dimm}) \), 
%		\( \lambda & (= \nmin \dimm) \),
		\( \lambda \) & Lagrange multiplier 		& 0, \textbf{0.05}, {0.1}, 0.5, 1 \\
		\( \Kern \) & the kernel & Epanechnikov: \( \Kern(t) = (1-t^{2})_{+} \) \\
		\( a \)	& decrease factor for \( \hiso_{k} \) & \( 2^{1/4}, \bm{2^{1/2}}, 2 \) \\
		\( k_{\max} \) & the number of repetitions in AO & \( \bm{3}, \, 5, \, 7 \) \\
		\( \hiso_{\min} \) & the smallest \( \hiso_{k} \) & \( s \sigmaX/\sqrt{n} \, \) for \( s=1, \, 2, \, \bm{3} \) \\
		\hline
    \end{tabular}
\caption{Tuning parameters for the single-index procedure}    
\label{TdefSI}
\end{table}

Step \( k \geq 0 \) of the procedure assumes the tensor \( \TEDR_{k} \) and the preliminary guess \( \betav_{\init} \)
available.

\begin{mylist}


\item
For each \( j \in \JJJ \), renew the sets of directions \( \Sph_{j} \);
%generate the sets of directions \( \Sph_{j} \) by drawing \( \sph \eqd \gaussv_{\TEDR_{k}}/\| \gaussv_{\TEDR_{k}} \| \)
%for \( \gaussv_{\TEDR_{k}} \sim \ND(0,\TEDR_{k}^{2}) \);

\item 
For each \( j \in \JJJ \), compute 
\( \neff_{j}^{(k)} \), \( \bigl( \Ima_{j,\sph}^{(k)} \), \( \Uv_{j,\sph}^{(k)} \bigr)_{\sph \in \Sph_{j}} \)
by \eqref{uyhvvy883hierwfewAVE} and \eqref{juu8hjfioee433erhAVE} with
\begin{EQA}[c]
	\weight_{ij} 
	= 
	\weight_{ij}^{(k)}
	= 
	\Kern\bigl( \| \TEDR_{k} (\Xv_{i} - \xv_{j}) \|^{2} \bigr) 
	\, .
\end{EQA}

\item 
With 
\( \Imav_{j}^{(k)} = \bigl( \Ima_{j,\sph}^{(k)} \bigr)_{\sph \in \Sph_{j}} \),
%\( \Sv_{j}^{(k)} = \bigl( S_{j,\sph}^{(k)} \bigr)_{\sph \in \Sph_{j}} \),
\( \UV_{j}^{(k)} = {(\Uv_{j,\sph}^{(k)})}_{\sph \in \Sph_{j}}^{\T} \),
solve
\begin{EQA}[c]
	\bigl\{ \betav_{k},(\fcl_{j}^{(k)})_{j \in \JJJ} \bigr\}
	=
	\arginf_{\betav \, , (\fclv_{j})_{j \in \JJJ}}
	\sum_{j \in \JJJ} \| \Imav_{j}^{(k)} - \fcl_{j} \, \UV_{j}^{(k)} \betav \bigr\|^{2} \neff_{j}^{(k)}
%	+ \frac{\lambda}{2} \| \betav - \betav_{k-1} \|^{2}
	\, 
	\qquad
\label{jkiiooofytkjhot43d}
\end{EQA}
using alternating optimization from Section~\ref{SLLaltern} which starts at \( \betav_{\init} \).
\item 
Compute the data fit
%\begin{EQA}[c]
%	\err_{k}
%	\eqdef
%	\sum_{j \in \JJJ} 
%	\bigl\{ Y_{i} - S_{j(i)}^{(k)} - \fcl_{j(i)}^{(k)} (\Xv_{i} - \Xbar_{j(i)}^{(k)})^{\T} \betav_{k} \bigr\}^{2} \, 
%	\, ,
%\label{jvfujkiiiut5tyerjfdk}
%\end{EQA}
\begin{EQA}[c]
	\err_{k}
	\eqdef
	\sum_{j \in \JJJ} 
	\bigl\{ Y_{i(j)} - S_{j}^{(k)} \bigr\}^{2} \, 
	\, ,
	\qquad
	S_{j}^{(k)}
	=
%	\biggl( \sumI \weight_{ij}^{(k)} \biggr)^{-1} 
	\frac{1}{\neff_{j}^{(k)}}
	\sumI Y_{i} \, \weight_{ij}^{(k)}
	\, .
	\qquad
\label{jvfujkiiiut5tyerjfdk}
\end{EQA}
%where
%\begin{EQA}[c]
%	\Xbar_{j}^{(k)}
%	=
%	\biggl( \sumI \weight_{ij}^{(k)} \biggr)^{-1} \sumI \Xv_{i} \, \weight_{ij}^{(k)}
%	\, ,
%	\qquad
%	S_{j}^{(k)}
%	=
%	\biggl( \sumI \weight_{ij}^{(k)} \biggr)^{-1} \sumI Y_{i} \, \weight_{ij}^{(k)}
%	\, .
%\end{EQA}

\item
If \( \hiso_{k}/a < \hiso_{\min} \) then terminate with output \( \betav_{\outp} = \betav_{k} \).

\item
Otherwise

\begin{mylist}
\item
Update \( \betav_{\init} = \betav_{k} \).

\item 
Increase \( k \) by one.
Update \( \hiso_{k} = \hiso_{k-1} / a \).

\item 
Redefine \( \TEDR_{k} \) in the form \eqref{iufhvu8kjbmnll475edf} with \( \betav = \betav_{\init} \) and \( \hiso = \hiso_{k} \):
\begin{EQA}[rcl]
	\TEDR_{k}^{2}
	&=&
	\hiso_{k}^{-2} \aniso_{k}^{2} \, \Id_{\dimd} 
	+ (1 - \aniso_{k}^{2}) \hiso_{k}^{-2} \betav_{\init} \, \betav_{\init}^{\T}
	\, ,
\label{iufhvu8kjbmnll475edfk}
\end{EQA}
where
\( \aniso_{k} \) is fixed by condition \eqref{jhgiverkytde4w34etycja}
with \( \TEDR = \TEDR_{k} \).
%\begin{EQA}[c]
%	\frac{1}{\nJJJ} \sum_{j \in \JJJ} \sumI \Kern\bigl( \| \TEDR_{k} (\Xv_{i} - \xv_{j}) \|^{2}
%	\bigr)
%	\geq 
%	\nmin
%	\, .
%	\qquad 
%\label{jhgiverkytde4w34etycja}
%\end{EQA}
\end{mylist}

\item
Loop.
% and \( (\fcn_{j}^{(k)},\fcl_{j}^{(k)})_{j \in \JJJ} \).
\end{mylist}


{\color{blue} 
\begin{mylist}
\item
Provide a possibility of choosing each tuning parameter (except \( \Kern \)).

\item
Run step-by-step, display the current values of \( \hiso_{k} \), \( \hiso_{k} /  \aniso_{k} \), \( \err_{k} \),
and \( |\cos(\betav_{k},\betavs)| \).


%\item
%Optimize the results w.r.t. the choice 
%of \( \nmin \), \( \nJJJ \), \( \nSph \), and \( \lambda \).

\item 
Plot the data \( Y_{i} \) against \( \Xv_{i}^{\T} \betav_{k} \) and against \( \Xv_{i}^{\T} \betavs \).

\end{mylist}
}

\Subsection{Stopping rule}
Our experimental study indicates that in some situations the quality of estimation 
degrades after a few steps.
In such cases, it would be reasonable to develop a data-driven stopping rule for selecting 
an optimal step \( k \) for terminating the procedure. 
One natural proposal is to define the final estimate \( \betav_{\outp} \) 
as \( \betav_{k} \) for the step \( k \) minimizing the data fit 
%\( \err_{\init} \) from \eqref{jvfujkiiiut5tyerjfdIni} and 
\( \err_{k} \) from \eqref{jvfujkiiiut5tyerjfdk}.
%Typically, we observe a unimodal behavior of \( \err_{k} \).
%This suggests a simplified procedure which terminates if \( \err_{k} \) starts to grow. 

\Subsection{Final results and comparison with other procedures}
{\color{blue} 
Compare the final results (initialization, full procedure, stopping rule) with 
\begin{mylist}
\item
Average Derivative Estimation (ADE);
\item
Sliced Inverse Regression (SIR);
\item
Minimum Average Variance Estimation (MAVE); Xia et al. (2002);
\end{mylist}
computed for the same data realization.
}

\Subsection{Breaking dimension and comparative study}

An important question for any dimension-reduction procedure is a breaking dimension.
Many such procedures exhibit reasonable performance for moderate values of \( \dimd \).
However, the results can degrade dramatically when the ambient dimension exceeds some value
depending on the case study. 

{\color{blue} 
For each setup, given by the parameters \( n \), \( \sigma_{\eps} \), \( \corr \),
run 100 simulations for each dimension \( \dimd = 10, 20, 30, 40, \ldots, 100 \).

Provide box plots of the final results 
after initialization, for the full procedure, and for the stopping rule.
Also, include the results for ADE, SIR, and MAVE. 
%\begin{mylist}
%\item
%Average Derivative Estimation (ADE).
%\item
%Sliced Inverse Regression (SIR).
%\item
%Minimum Average Variance Estimation (MAVE); Xia et al. (2002).
%\end{mylist}
%
%
}



%\Subsection{Data fit and stopping rule}

%\end{document}
%%%%%%%%%%%%%%%%%%%%%%%%%%%%%%%  new section  %%%%%%%%%%%%%%%%%%%%%%%%%%%%%%%
\Section{Multi-index model}
\label{Smultind}


Consider the e.d.r. regression problem
\begin{EQA}
	Y_{i}
	&=&
	\gs(\Xv_{i}) + \eps_{i}
	=
	\fs(\PEDR \Xv_{i}) + \eps_{i}
	\, ,
\label{YifXieiedr}
\end{EQA}
where \( \Xv_{i} \in \R^{\dimd} \) for large \( \dimd \), 
\( \PEDR \) is an unknown element from the set \( \PROJ_{\dimd,\dimm} \) 
of orthogonal mappings from 
from \( \R^{\dimd} \) to \( \R^{\dimm} \) for \( \dimm \ll \dimd \), and
\( \fs \) is a function on \( \R^{\dimm} \).
The mapping \( \PEDR \) can be written as 
\begin{EQA}
	\PEDR \xv
	&=&
	(\betav_{1}^{\T} \xv) \betav_{1} + \ldots + (\betav_{\dimm}^{\T} \xv) \betav_{\dimm}
	\, ,
\label{ydgenb8jejohkou90tkr}
\end{EQA}
where \( \betav_{m}^{\T} \) are the orthonormal rows of \( \PEDR \).
Here the identifiability issue is even more severe than in the single-index model, because 
the set of the vectors \( \betav_{j} \) can be exchanged, rotated, etc. 

Similarly to the single-index case, 
the structural information about the model
can be recovered from the gradient vectors \( \nabla \gs(\xv) \).
More precisely, under \eqref{YifXieiedr}, the \( \dimd \times \dimd \)-matrix
\begin{EQA}[c]
	\JEDR
	\eqdef
	\sum_{j \in \JJJ} \nabla \gs(\xv_{j}) \, \nabla \gs(\xv_{j})^{\T} \, \neff_{j}
	=
	\sum_{j \in \JJJ} \PEDR^{\T} \fclv_{j} \, \bigl( \PEDR^{\T} \fclv_{j} \bigr)^{\T} \, \neff_{j}
%	\biggl( \sum_{m=1}^{\dimm} \fcl_{j,m} \, \betav_{m} \biggr)^{\T}
	=
	\PEDR^{\T} \biggl( \sum_{j \in \JJJ} \fclv_{j} \, \fclv_{j}^{\T}  \, \neff_{j} \biggr) \PEDR
	\, 
	\qquad
\label{ihjfiojihfer6t4rrt}
\end{EQA}
with \( \fclv_{j} = \nabla \fs(\PEDR \xv_{j}) \)
describes the multi-index structure.
Moreover, 
if all positive eigenvalues of \( \JEDR \) are different, 
then the SVD \( \JEDR = {\PEDRs}^{\T} \Lambda \, \PEDRs \) with \( \Lambda = \diag(\lambda_{1} \geq \ldots \geq \lambda_{\dimm}) \)
provides a unique canonical representation for the multi-index model \eqref{YifXieiedr}.


%%%%%%%%%%%%%%%%%%%% section %%%%%%%%%%%%%%%%%%%%
\Subsection{Preliminaries}
For a trainign set \( \II \), a collection of centers \( \{ \xv_{j} \}_{j \in \JJJ} \),
and a localizing kernel \( \Kern \), consider the localizing weights of the form
\begin{EQA}
	\weight_{ij}
	&=&
	\Kern\bigl( \| \TEDR(\Xv_{i} - \xv_{j}) \|^{2} \bigr) 
	\, ,
	\qquad
	i \in \II,
	\;
	j \in \JJJ
	\, ,
\label{s7cuqkdwywdhbjwbfjwbdjm}
\end{EQA}
where \( \TEDR \) is a localizing tensor \( \TEDR \) (symmetric \( \dimd \)-matrix)
which reflects the available information about the multi-index structure in \eqref{YifXieiedr}.
At the beginning, an uninformative isotropic tensor \( \TEDR_{0} = \hiso_{0}^{-1} \Id_{\dimd} \) will be used.
During iteration, we learn the structure and use it to improve \( \TEDR \); see Section~\ref{SEDRsa} ahead.

For every \( j \in \JJJ \), assume a set \( \Sph_{j} \) or directions \( \sph \) (unit vectors in \( \R^{\dimd} \)) 
to be fixed.
Compute for all \( j \in \JJJ \), \( \sph \in \Sph_{j} \),
the local means \( \Xbar_{j} \) from \eqref{uyhvvy883hierwfewAVE} and 
the quantities \( \Ima_{j,\sph}, S_{j,\sph} \), and \( \Uv_{j,\sph} \) 
from \eqref{juu8hjfioee433erhAVE} as in the single-index case.
%\begin{EQ}[rcl]
%	\Ima_{j,\sph}
%	&=&
%	\sumI Y_{i} \langle \Xv_{i} - \xv_{j} \, ,\sph \rangle \, \weight_{ij} 
%	\, ,
%	\qquad
%	S_{j,\sph}
%	=
%	\sumI \langle \Xv_{i} - \xv_{j} \, ,\sph \rangle \, \weight_{ij}
%	\, ,
%%\label{ytd7e7ydwuhd5eythdYm}
%	\\
%	\Uv_{j,\sph}
%	&=&
%	\sumI (\Xv_{i} - \xv_{j}) \, \langle \Xv_{i} - \xv_{j} \, ,\sph \rangle \, \weight_{ij}
%	\in \R^{\dimd}
%	\, .
%	\qquad
%\label{ytd7e7ydwuhd5eythdm}
%\end{EQ}
Model equation \eqref{YifXieiedr} implies
\begin{EQA}
	\Ima_{j,\sph}
	&=&
	\frac{1}{\neff_{j}}
	\sumI Y_{i} \langle \Xv_{i} - \Xbar_{j} \, ,\sph \rangle \, \weight_{ij}  
	\\
	&=&
	\frac{1}{\neff_{j}}
	\sumI \fs(\PEDR \Xv_{i}) \langle \Xv_{i} - \Xbar_{j} \, ,\sph \rangle \, \weight_{ij}
	+ 
	\frac{1}{\neff_{j}}
	\sumI \eps_{i} \langle \Xv_{i} - \Xbar_{j} \, ,\sph \rangle \, \weight_{ij} \, .
\label{hgghjkklvddfrdshweym}
\end{EQA}
%
%as in \eqref{ytd7e7ydwuhd5eythdm}.
The function \( \fs \) is now \( \dimm \)-variate, and we consider a linear approximation 
\begin{EQA}
	\fs(\PEDR \xv) 
	& \approx &
	\fcn_{j} + \fclv_{j}^{\T} \PEDR (\xv - \xv_{j})
	=
	\fcn_{j} + \biggl( \sum_{m=1}^{\dimm} \fcl_{j,m} \, \betav_{m} \biggr)^{\T} (\xv - \xv_{j}) ,
\label{ucjcvjuvyvyf6r6rwjdy}
\end{EQA}
where \( \fclv_{j} = (\fcl_{j,m}) \in \R^{\dimm} \).
Correspondingly, with \( \Xbar_{j} \) from \eqref{uyhvvy883hierwfewAVE},
\begin{EQA}[rcl]
	&& \nquad
	\frac{1}{\neff_{j}}
	\sumI \fs(\PEDR \Xv_{i}) \, \langle \Xv_{i} - \Xbar_{j} \, ,\sph \rangle \, \weight_{ij} 
	\approx 
	\frac{1}{\neff_{j}}
	\sumI \bigl\{ \fcn_{j} 
%	+ \sum_{m=1}^{\dimm} \fcl_{j,m} \, (\Xv_{i} - \xv_{j})^{\T} \betav_{m} \Bigr\}
	+ \fclv_{j}^{\T} \PEDR (\Xv_{i} - \xv_{j})
	\langle \Xv_{i} - \Xbar_{j} \, ,\sph \rangle \bigr\} \weight_{ij} 
	\\
	&=&
	\fclv_{j}^{\T} \PEDR \, \Uv_{j,\sph}
	=
%	\fcn_{j} S_{j,\sph} + 
	\Uv_{j,\sph}^{\T} \PEDR^{\T} \, \fclv_{j}
	\, .
\label{yd6x5ewbduwuqf64wfgfymi}
\end{EQA}
%where \( \fclv_{j} \cdot \betav = \fcl_{j,1} \, \betav_{1} + \ldots + \fcl_{j,\dimm} \, \betav_{\dimm} \).
The full dimensional parameter \( \prmtv \) includes
%the corresponding loading factors \( \mue_{1},\ldots,\mue_{\dimm} \), 
\( (\imav_{j,\sph}) \in \R^{\nJJJ \nSph} \), 
the coefficients \( (\fclv_{j})_{j \in \JJJ} \in \R^{\nJJJ\dimm} \),
and the target mapping \( \PEDR = (\betav_{1},\ldots,\betav_{\dimm})^{\T} \).
Denote \( \imav_{j} = \bigl( \ima_{j,\sph} \bigr)_{\sph \in \Sph_{j}} \in \R^{\nSph} \), 
\( \Imav_{j} = \bigl( \Ima_{j,\sph} \bigr)_{\sph \in \Sph_{j}} \in \R^{\nSph} \),
%\( \Sv_{j} = \bigl( S_{j,\sph} \bigr)_{\sph \in \Sph_{j}} \in \R^{\nSph} \),
and let \( \UV_{j} \) be the \( \nSph \times \dimd \)-matrix with the rows \( \Uv_{j,\sph}^{\T} \),
\( \sph \in \Sph_{j} \).
Then the objective function reads
\begin{EQA}
	\LL(\prmtv)
	&=&
	\frac{1}{2} \sum_{j \in \JJJ} \Bigl( \| \Imav_{j} - \imav_{j} \|^{2} 
		+ \bigl\| \imav_{j} - \UV_{j} \, \PEDR^{\T} \fclv_{j} \bigr\|^{2} \Bigr) \neff_{j}
	\, .
	\qquad
%\label{gtgyuy6ew5dtde32322jtmi}
\end{EQA}
The auxillary variables \( \zv_{j} \) can be eliminated, leading to
\begin{EQA}[rcl]
	\LL(\prmtv)
	&=&
	\LL(\PEDR,(\fclv_{j})_{j \in \JJJ})
	=
	\frac{1}{4} \sum_{j \in \JJJ} \bigl\| \Imav_{j} - \UV_{j} \, \PEDR^{\T} \fclv_{j} \bigr\|^{2} \neff_{j}
	\, .
	\qquad
\label{gtgyuy6ew5dtde32322jtmiG0}	
\end{EQA}
This function should be minimized over the set of all coefficients \( (\fclv_{j})_{j \in \JJJ} \) and all
orthonormal mappings \( \PEDR \in \PROJ_{\dimd,\dimm} \) with \( \PEDR \PEDR^{\T} = \Id_{\dimm} \):
\begin{EQA}[c]
%	\prmtv_{0}
%	= 
%	\{ \PEDR_{0} \, , (\fcn_{j}^{(0)},\fclv_{j}^{(0)})_{j \in \JJJ} \} 
%	=
	\LL(\prmtv)
	\to
	\inf_{\PEDR \, , (\fclv_{j})_{j \in \JJJ}} 
	\, .
\label{ju7fudeju74334fhfd}
\end{EQA}

%%%%%%%%%%%%%%%%%%%% section %%%%%%%%%%%%%%%%%%%%
\Subsection{Alternating optimization}
\label{SaltMI}
Solving \eqref{ju7fudeju74334fhfd} is challenging because
the class of mappings \( \PROJ_{\dimd,\dimm} \) is not convex; this makes the problem computationally hard.
We also face an identifiability issue: a solution \( \PEDR \) corresponds to 
a particular choice of the basis in the EDR space \( \EDR \).
This basis can be arbitrarily rotated.
We apply the alternating procedure similar to Section~\ref{SLLaltern}.
The basic idea of the proposed approach is to relax the target \( \PEDR \in \PROJ_{\dimd,\dimm} \)
and optimize \( \LL(\prmtv) \) w.r.t. 
\( \BEDR = (\betav_{1},\ldots,\betav_{\dimm})^{\T} \in \R^{\dimm \times \dimd} \). 
Then the desired orthogonal mapping \( \PEDR \) will be obtained using SVD 
of the solution \( \BEDR \).
%
%To ensure numerical stability, we add a penalty term 
%\( \lambda \| \BEDR - \PEDR_{\init} \|_{\Fr}^{2} \),
%where \( \PEDR_{\init} \) is either a starting guess or 
%a previous step solution \( \PEDR_{\init} \).

The alternating optimization (AO) procedure assumes 
a starting guess \( \PEDR_{\init} \) to be given.
It repeats several times (\( k_{\max} \)) one-step alternating, which reads as follows.

\paragraph{One-step alternating} 

\begin{mylist}
\item
With \( \PEDR_{\init} \) fixed, compute for every \( j \in \JJJ \)
\begin{EQA}[c]
	\hat{\fclv}_{j}
	=
	\arginf_{\fclv_{j}} \bigl\| \Imav_{j} - \UV_{j} \PEDR_{\init}^{\T} \fclv_{j} \bigr\|^{2}
	\, .
\end{EQA}


\item 
Then, with \( (\hat{\fclv}_{j})_{j \in \JJJ} \) fixed, compute \( \hat{\BEDR} \):
\begin{EQA}[c]
	\hat{\BEDR}
	=
	\arginf_{\BEDR \in \R^{\dimm \times \dimd}} 
		\biggl\{ \sum_{j \in \JJJ} \bigl\| \Imav_{j} - \UV_{j} \BEDR^{\T} \hat{\fclv}_{j} \bigr\|^{2} \neff_{j}
	+ \lambda \| \BEDR - \PEDR_{\init} \|_{\Fr}^{2} 
	\biggr\}	
	\, ,
\end{EQA}


\item
Compute 
\begin{EQA}[rcl]
	\JEDR
	\eqdef
	\sum_{j \in \JJJ} (\hat{\BEDR}^{\T} \, \hat{\fclv}_{j}) \, \bigl( \hat{\BEDR}^{\T} \, \hat{\fclv}_{j} \bigr)^{\T} \neff_{j}
	&=&
	\hat{\BEDR}^{\T} \Bigl( \sum_{j \in \JJJ} \hat{\fclv}_{j} \, \hat{\fclv}_{j}^{\T} \neff_{j} \Bigr) \hat{\BEDR}
\end{EQA}
and its SVD
\begin{EQA}[c]
	\JEDR
	=
	\PEDR^{\T} \, \Lambda \, \PEDR
	\, ,
	\qquad
	\PEDR \, \PEDR^{\T}
	= 
	\Id_{\dimm}
	\, ,
	\qquad
	\Lambda = \diag(\lambda_{1} \geq \ldots \geq \lambda_{\dimm})
	\, .
\end{EQA}

\item 
Output \( \PEDR_{\outp} = \PEDR \) and \( \Lambda_{\outp} = \Lambda/\| \Lambda \| = \lambda_{1}^{-1} \Lambda \).
\end{mylist}


\paragraph{Alternating optimization procedure} repeats
one-step alternating several times.
At every iteration, the output \( \PEDR_{\outp} \) of the previous step is used as input \( \PEDR_{\init} \).
The procedure delivers the last step results \( \PEDR_{\outp} \) and \( \Lambda_{\outp} \) after \( k_{\max} \) iterations.

%{\color{blue} check one-step improvement and final quality of \( \PEDR \)}


%%%%%%%%%%%%%%%%%%%% section %%%%%%%%%%%%%%%%%%%%
\Subsection{Quality of estimation}
\label{SEDRqua}
This section discusses how the discrepancy \( \epsilon(\PEDR , \PEDRs) \) between the true projector \( \PEDRs \) 
from \eqref{ihjfiojihfer6t4rrt} and 
the estimated projector \( \PEDR \) can be defined.
%The value \( \epsilon(\PEDR , \PEDR_{\jl}) \) measures .
The Frobenius or operator norms of the difference \( \PEDR - \PEDRs \) cannot be used here because
\( \PEDR \) is not uniquely defined; one can swap the sign of each row without changing the manifold projection. 
This suggests
\begin{EQA}[rcl]
	\epsilon(\PEDR , \PEDRs)
	& \eqdef &
	\bigl\| \PEDR (\Id_{\dimd} - {\PEDRs}^{\T} \PEDRs) \bigr\|_{\Fr}^{2}
	=
	\tr \bigl\{ (\Id_{\dimd} - {\PEDRs}^{\T} \PEDRs) \PEDR^{\T} \PEDR (\Id_{\dimd} - {\PEDRs}^{\T} \PEDRs) \bigr\}
	\\
	&=&
	\tr \bigl\{ \PEDR^{\T} \PEDR (\Id_{\dimd} - {\PEDRs}^{\T} \PEDRs) \bigr\}
	\, .
\label{hjjjuuuu55gswwdtgf}
\end{EQA}

\Subsection{Data fit}
The proposed procedure is iterative and involves several steps.
As in the single-index case, 
a proper stopping rule is an issue and requires evaluating the quality of estimation at every step.
%
The quality of estimation will be measured by the generalization error on the testing set
which corresponds to the set \( \JJJ \) of centers \( \xv_{j} \).
%
For each \( j \in \JJJ \), assume that the point \( \Xv_{i(j)} \) is aligned  
with a center \( \xv_{j} \).
Then the data fit \( \err \) from \eqref{jvfujkiiiut5tyerjfd} can be used here as well.
%For the initialization step, we apply \eqref{jvfujkiiiut5tyerjfdIni} and \eqref{ihvifvityguwhd8354e}.

%%%%%%%%%%%%%%%%%%%% section %%%%%%%%%%%%%%%%%%%%
\Subsection{Structural adaptation}
\label{SEDRsa}
The use of isotropic weights \( \weight_{ij} \) for computing the observables 
\( \Imav_{j} \), \( \Sv_{j} \), and \( \UV_{j} \) is suboptimal and leads to poor 
quality of recovering the EDR-space \( \EDR \).
The structure-adaptive approach from \cite{HJPS} suggests extracting the structural information from the current-step estimate 
and using it for improving the accuracy of the next-step estimate by applying 
anisotropic weights \( \weight_{ij} = \Kern(\| \TEDR (\Xv_{i} - \xv_{j}) \|^{2}) \).
The tensor \( \TEDR \) here accumulates the structural information learnt at previous steps. 
At the beginning, we use \( \TEDR = \hiso_{0}^{-1} \Id_{\dimd} \).
Let now \( k-1 \) steps of the procedure have been performed.
The results are given as matrices \( \PEDR_{k-1} \) and \( \Lambda_{k-1} \) with \( \| \Lambda_{k-1} \| = 1 \).
For the step \( k \),
with \( \JEDR_{k-1} = \PEDR_{k-1}^{\T} \Lambda_{k-1} \PEDR_{k-1} \) and \( \aniso_{k} < 1 \), define 
\begin{EQA}[c]
	\TEDR_{k}^{2}
	\eqdef
	\aniso_{k}^{2} \, \hiso_{k}^{-2} \Id_{\dimd} + \hiso_{k}^{-2} \, \JEDR_{k-1}
	\quad
	\text{or }
	\quad
	\TEDR_{k}^{2}
	\eqdef
	\aniso_{k}^{2} \, \hiso_{k}^{-2} (\Id_{\dimd} - \PEDR_{k-1}^{\T} \PEDR_{k-1}) + \hiso_{k}^{-2} \, \JEDR_{k-1}
	\, .
\end{EQA}
For any \( \uv \in \R^{\dimd} \), this means
\begin{EQA}[rcl]
	\| \TEDR_{k} \uv \|^{2}
	&=&
	\aniso_{k}^{2} \, \hiso_{k}^{-2} \, \| \uv \|^{2} + \hiso_{k}^{-2} \, \| \Lambda_{k-1}^{1/2} \PEDR_{k-1} \uv \|^{2}
	\\
	\text{or }
	\quad
	\| \TEDR_{k} \uv \|^{2}
	&=&
	\aniso_{k}^{2} \, \hiso_{k}^{-2} \, \bigl( \| \uv \|^{2} - \| \PEDR_{k-1} \uv \|^{2} \bigr) 
	+ \hiso_{k}^{-2} \, \| \Lambda_{k-1}^{1/2} \PEDR_{k-1} \uv \|^{2}
	\, .
\end{EQA}
%
%These matrices can be used to improve the procedure by redefining the tensor \( \TEDR \) as
%\begin{EQA}
%	\TEDR^{2}
%	&=&
%	\hiso^{-2} \bigl( \aniso^{2} \, \Id_{\dimd} + \JEDR_{0} \bigr)
%%	\sum_{m} \biggl( \sum_{j \in \JJJ} \fcl_{j,m}^{2} \biggr) \betav_{m} \betav_{m}^{\T} 
%	\, ,
%\label{yfijowjdo76dvuienk}
%\end{EQA}
%This value describes the degree of kernel anisotropy.
%Alternatively, %one can use
%\begin{EQA}[c]
%	\TEDR^{2}
%	=
%	\hiso^{-2} \aniso^{2} (\Id_{\dimd} - \PEDR_{0}^{\T} \PEDR_{0}) + \hiso^{-2} \, \JEDR_{0}
%	\, .
%\end{EQA}
This tensor leads to anisotropic weights 
\( \weight_{ij}^{(k)} = \Kern\bigl( \| \TEDR_{k}(\Xv_{i} - \xv_{j}) \|^{2} \bigr) \)
which strongly penalize in the principal directions of the estimated EDR space.
%It will also be used for renewal of the set of directions \( \Sph_{j} \) which are randomly sampled as
%\( \gaussv_{\TEDR_{k}}/\| \gaussv_{\TEDR_{k}} \| \) for \( \gaussv_{\TEDR_{k}} \sim \ND(0,\TEDR_{k}^{2}) \). 
The sums \( \Ima_{j,\sph} \) and \( \Uv_{j,\sph} \) in \eqref{juu8hjfioee433erhAVE}
should be recomputed with these changes.
The improved estimator \( \prmtv_{k} \) is given by 
minimization of \( \LL(\prmtv) \) from \eqref{gtgyuy6ew5dtde32322jtmiG0}
after this update.
%
This structural adaptation step can be iterated several times.
The values \( \hiso_{k}, \aniso_{k} \) may change during iterations: 
\( \hiso_{k} \) decreases at every step with a fixed factor 
\( a > 1 \), while \( \aniso_{k} \) is fixed by condition \eqref{jhgiverkytde4w34etycja}.

%%%%%%%%%%%%%%%%%%%% section %%%%%%%%%%%%%%%%%%%%
\Subsection{The procedure}
\label{SEDRmiproc}

We consider the same design as in the single-index case; see Section~\ref{SdesignSI}.
Only the function \( \fs \) has to be \( \dimm \)-variate;
e.g. additive \( \fs(x_{1},x_{2}) = x_{1}^{2} + \sin(s x_{2}) \)
or multiplicative \( \fs(x_{1},x_{2}) = x_{1} \sin(s x_{2}) \) for \( s=1,2,3,4 \).
Instead of the vectors \( \betavs \) we generate randomly a collection 
\( \betavs_{1},\ldots,\betavs_{\dimm} \). 

Initialization is an important issue.
A natural proposal is to apply local linear gradient estimation
from Section~\ref{Siniloclin} and define 
\( \PEDR_{\init} \) using the first \( \dimm \) principal directions
for the set of the estimated gradients \( \hat{\nabla \gs}(\xv_{j}) \), \( j \in \JJJ \).




\paragraph{Initialization (requires to fix \( \nJJJ, \nlin, \displ \)); see Table~\ref{TdefSIini}}
\begin{mylist}
\item
Generate \( (\xv_{j} \, , j \in \JJJ) \).

\item
Fix \( \PEDR_{\init} \) by \( \dimm \)-dimensional PCA on the set of local gradient estimates
\( \hat{\nabla \gs}(\xv_{j}) \) from \eqref{ihvifvityguwhd8354e}.
%; see Section~\ref{Siniloclin}.

%\item
%Compute \( \err_{\init} \) by \eqref{jvfujkiiiut5tyerjfdIni}.

\item
Fix \( \hiso_{0} \) as the smallest value ensuring \eqref{jhgiverkytde4w34etycj}.
Define \( \TEDR_{0} = \hiso_{0}^{-1} \Id_{\dimd} \).

\item Set \( k=0 \).

\end{mylist}

{\color{blue} 
\begin{mylist}
\item 
Display the first \( \dimm + 1 \) eigenvalues of \( \JEDR \) from \eqref{uhcudiiy66yhjikfed}.

\item 
Display the error \( \epsilon(\PEDR , \PEDRs) \) from \eqref{hjjjuuuu55gswwdtgf}.

\item
Provide the possibility of manual selection of \( \nlin \) and of \( \nJJJ \).
\end{mylist}
}

%%%%%%%%%%%%%%%%%%%% section %%%%%%%%%%%%%%%%%%%%
\Subsubsection{Main procedure}
\label{SEDRmiproc}


Fix the set of tuning parameters; see Table~\ref{TdefMI}.
\begin{table}[h]
\begin{tabular}{c|c|c}
        \hline
        Parameter & Meaning & possible values and default \\
        \hline
        \( \nmin \) & averaged mass of a local vicinity & \( 7,10, 15, \bm{20} \) \\
        \( \nJJJ \) & size of \( \JJJ \) 				& \( s \, n/\nmin \) for \( s=1, \ldots, \bm{\nmin/2}, \ldots, \nmin \) \\
		\( \nSph \) & size of each \( \Sph_{j} \)		& \( \nmin, \bm{2\nmin}, 3\nmin \) \\
%\( a (= 2^{1/\dimm}) \), 
%		\( \lambda & (= \nmin \dimm) \),
		\( \lambda \) & Lagrange multiplier 		& 0, \textbf{0.05}, {0.1}, 0.5, 1 \\
		\( \Kern \) & the kernel & Epanechnikov: \( \Kern(t) = (1-t^{2})_{+} \) \\
		\( k_{\max} \) & the number of repetitions in AO & \( 3, \bm{5}, 7 \) \\
		\( a \)	& decrease factor for \( \hiso_{k} \) & \( 2^{1/(\dimm+2)}, {2^{1/(\dimm+1)}}, \bm{2^{1/\dimm}} \) \\
		\( \hiso_{\min} \) & the smallest \( \hiso_{k} \) & \( s \sigmaX/\sqrt{n} \, \) for \( s=1,2,\bb{3} \) \\
		\hline
    \end{tabular}
\caption{Tuning parameters for the multi-index procedure}    
\label{TdefMI}
\end{table}

%\( \dimm (\leq 3) \), 
%\( \nmin (= 10) \), 
%\( \nJJJ (= n) \), 
%\( \nSph (= \min\{ \nmin , \dimd \}) \),
%\( a (= 2^{1/\dimm}) \), 
%\( \lambda (= \dimm) \),
%\( k_{\max} (= 5) \).

\paragraph{Step \( k \geq 0 \) (tensor \( \TEDR_{k} \) and the preliminary guess \( \PEDR_{\init} \) are assumed 
%\( \PEDR_{k-1} \) and \( \Lambda_{k-1} \) 
available.)}

\begin{mylist}


\item
For all \( j \in \JJJ \), 
redraw directions \( \sph \in \Sph_{j} \) with
\( \sph \eqd \gaussv/\| \gaussv \| \), \( \gaussv \sim \ND(0,\Id_{\dimd}) \).
%\( \sph \eqd \gaussv_{\TEDR_{k}}/\| \gaussv_{\TEDR_{k}} \| \), \( \gaussv_{\TEDR_{k}} \sim \ND(0,\TEDR_{k}^{2}) \).

\item 
For each \( j \in \JJJ \), compute 
\( \neff_{j}^{(k)} \), \( \bigl( \Ima_{j,\sph}^{(k)} \), \( \Uv_{j,\sph}^{(k)} \bigr)_{\sph \in \Sph_{j}} \)
by \eqref{uyhvvy883hierwfewAVE} and \eqref{juu8hjfioee433erhAVE} with
\begin{EQA}[c]
	\weight_{ij} 
	= 
	\weight_{ij}^{(k)}
	= 
	\Kern\bigl( \| \TEDR_{k} (\Xv_{i} - \xv_{j}) \|^{2} \bigr) 
	\, .
\end{EQA}

\item 
With 
\( \Imav_{j}^{(k)} = \bigl( \Ima_{j,\sph}^{(k)} \bigr)_{\sph \in \Sph_{j}} \),
%\( \Sv_{j}^{(k)} = \bigl( S_{j,\sph}^{(k)} \bigr)_{\sph \in \Sph_{j}} \),
\( \UV_{j}^{(k)} = {(\Uv_{j,\sph}^{(k)})}_{\sph \in \Sph_{j}}^{\T} \),
solve
\begin{EQA}[c]
	\bigl\{ \PEDR_{k},(\fclv_{j}^{(k)})_{j \in \JJJ} \bigr\}
	=
	\arginf_{\PEDR,(\fclv_{j})_{j \in \JJJ}}
	\sum_{j \in \JJJ} \bigl\| \Imav_{j}^{(k)} - \UV_{j}^{(k)} \, \PEDR^{\T} \fclv_{j} \bigr\|^{2} \neff_{j}^{(k)}
%	+ \lambda \| \BEDR - \PEDR_{k-1} \|_{\Fr}^{2}
	\, 
\end{EQA}
%where \( \PEDR %= (\betav_{1}, \ldots, \betav_{\dimm})^{\T} \in \R^{\dimm \times \dimd} \), 
%\in \PROJ_{\dimd, \dimm} \), 
%\( \fcn_{j} \in \R \), \( \fclv_{j} \in \R^{\dimm} \), \( j \in \JJJ \),
using the alternating optimization procedure as in Section~\ref{SaltMI} starting from 
\( \PEDR_{\init} \).
This step delivers \( \PEDR_{k} \) and \( \Lambda_{k} \) with \( \| \Lambda_{k} \| = 1 \).

\item 
Compute the data fit
\begin{EQA}[c]
	\err_{k}
	\eqdef
	\sum_{j \in \JJJ} \bigl| Y_{i(j)} - S_{j}^{(k)} \bigr|^{2}
%	\bigl\{ Y_{i} - S_{j}^{(k)} - (\Xv_{i} - \Xbar_{j}^{(k)})^{\T} \PEDR_{k}^{\T} \fclv_{j}^{(k)} \bigr\}^{2} \, 
%	\weight_{ij}^{(k)}
	\, ,
\label{jvfujkiiiut5tyerjfdkmi}
%\end{EQA}
%where
%\begin{EQA}[c]
%	\Xbar_{j}^{(k)}
%	=
%	\biggl( \sumI \weight_{ij}^{(k)} \biggr)^{-1} \sumI \Xv_{i} \, \weight_{ij}^{(k)}
%	\, ,
	\qquad
	S_{j}^{(k)}
	=
	\frac{1}{\neff_{j}^{(k)}} \sumI Y_{i} \, \weight_{ij}^{(k)}
	\, .
\end{EQA}

\item
Terminate if \( \hiso_{k}/a < \hiso_{\min} \).
Output \( \PEDR_{k} \), \( \Lambda_{k} \).
%, \( (\fcn_{j}^{(k)},\fclv_{j}^{(k)})_{j \in \JJJ} \).

\item 
%:
Otherwise:
\begin{mylist}


\item
Update \( \PEDR_{\init} = \PEDR_{k} \).

\item
Increase \( k \) by one and define \( \hiso_{k} = \hiso_{k-1} / a \).


\item
%With \( \JEDR_{k-1} = \PEDR_{k-1}^{\T} \Lambda_{k-1} \PEDR_{k-1} \), 
Define \( \TEDR_{k} \) by
\begin{EQA}[c]
	\TEDR_{k}^{2}
	\eqdef
%	\aniso_{k}^{2} \, \hiso_{k}^{-2} \Id_{\dimd} + \hiso_{k}^{-2} \, \JEDR_{k-1}
%	\quad
%	\text{or }
%	\quad
	\aniso_{k}^{2} \, \hiso_{k}^{-2} (\Id_{\dimd} - \PEDR_{k-1}^{\T} \PEDR_{k-1}) 
	+ \hiso_{k}^{-2} \, \PEDR_{k-1}^{\T} \Lambda_{k-1} \PEDR_{k-1}
	\, .
\end{EQA}
where \( \aniso_{k} \) is the smallest value ensuring \eqref{jhgiverkytde4w34etycja}.
%\begin{EQA}[c]
%	\frac{1}{\nJJJ} \sum_{j \in \JJJ} \sumI 
%		\Kern\bigl( \| \TEDR_{k} (\Xv_{i} - \xv_{j}) \|^{2} \bigr)
%%	=
%%	\frac{1}{\nJJJ}\sum_{j \in \JJJ} \sumI \weight_{ij}
%	\geq 
%	\nmin
%	\, .
%\label{jhgiverkytde4w34etycjam}
%\end{EQA}
\end{mylist}
\item Loop.
\end{mylist}

{\color{blue} 
\begin{mylist}
\item
Provide a possibility of manually choosing each tuning parameter (except \( \Kern \)).

\item
Run step-by-step, display the current values of \( \hiso_{k} \), \( \hiso_{k} /  \aniso_{k} \),
the accuracy \( |\epsilon(\PEDR_{k},\PEDRs)| \), and the fit \( \err_{k} \).


\item
Provide the stopping diagnostics as in the single-index case.

\item
Compare the final results (initialization, full procedure, stopping rule) with 
%Average Derivative Estimation (ADE);
SIR and MAVE computed for the same data realization.

%Enable optimal selection
%of \( \nmin \), \( \nJJJ \), \( \nSph \), and \( \lambda \).
%

\end{mylist}

}
